\begin{abstract}
\noindent Does the yield advantage of irrigated maize grow with extreme summer heat? French agricultural statistics report grain maize yields \emph{separately} for irrigated and rainfed land, by department. I relate the gap between the two, 2000--2025, to degree-days above 32\,$^{\circ}$C in July--August measured on the parcels where maize grows. The specification was written after an exploratory analysis of 2010--2025 and before any estimate on 2000--2009. In 17 departments outside the South-West (\QoneN{} department-years), each degree-day is associated with a \QonePct\,\% higher ratio of irrigated to rainfed yield, so that the median advantage of \GapMedianPct\,\% would become \GapMedianAfterPct\,\% (one-sided wild cluster bootstrap by year, $p = \QoneP$; measured size \QoneSizeWCR\,\%; 90\,\% interval \CiNinetyLow{} to \CiNinetyHigh\,\%). This meets the pre-specified rule, with four reservations. \ExploSharePct\,\% of the sample comes from the exploratory years, and the new decade keeps the sign at about half the estimate ($p = \QoneReplP$). Without the six departments with uplands the result is not significant; a cluster-jackknife statistic is not significant against a Student distribution ($p = \CvThreeP$) but is when bootstrapped ($p = \CvBootP$, size \CvBootSize\,\%). In the older release of the statistics, the irrigated-to-rainfed ratio repeats from one year to the next more often than chance, as an imputation would make it, and a pre-declared check classes this as a live threat. And with the vapour pressure deficit in the model, the association moves to the drying power of the air: it concerns hot, dry summers rather than heat alone. It holds with station-based temperatures and with satellite surface temperature, replicates on silage maize, and all \MultiN{} specifications of a multiverse analysis give a positive estimate; early, coarse estimates for the hot summer of 2026 point the same way. A pre-specified interaction with water restrictions has the opposite sign to the hypothesis. The result is an association, not the causal effect of irrigation. A methodological check changed the conclusions: standard errors clustered by department, or Conley, reject a true null in \UsualSizeMin\,\% to \UsualSizeMax\,\% of placebo simulations on these panels, and with them heat would have appeared ``significantly'' associated with lower wheat yields ($p = \WheatPdep$, against $p = \WheatPwcr$ with a correctly sized test). Wheat yields over 1959--2025 and out-of-sample forecasts give null results, reported in full.

\medskip\noindent\textbf{Keywords:} heat stress; irrigation; maize; France; panel data; wild cluster bootstrap; pre-specification.\\
\textbf{JEL codes:} Q15, Q54, C23, C12.
\end{abstract}
